\section{Some types of spaces having the property (GDF)}

We already know that every Fréchet space has the property (GDF) (TVS, I, §3, No.~3, Cor. 5 of Th. 1).

PROPOSITION 2. --- Let E be a vector space, \((\mathrm{F}_{\alpha})_{\alpha\in\mathrm{A}}\) a family of locally compact spaces having the property (GDF), and for each \(\alpha\in A\) let \(h_{\alpha}\) be a linear mapping of \(F_{\alpha}\) into E. If E is equipped with the finest locally convex topology that makes the \(h_{\alpha}\) continuous, then E has the property (GDF).

Let \(u\) be a linear mapping of \(\mathbf{E}\) into a Banach space \(\mathbf{B}\) , such that every limit in \(\mathbf{E} \times \mathbf{B}\) of every convergent sequence of points of the graph \(\Gamma\) of \(u\) also belongs to \(\Gamma\) . It suffices to show that, for every \(\alpha \in \mathbf{A}\) , \(u \circ h_{\alpha}\) is continuous on \(\mathrm{F}_{\alpha}\) (TVS, II, §4, No.~4, Prop. 5). Now, let \((x_n)\) be a sequence of elements of \(\mathrm{F}_{\alpha}\) having a limit \(a\) and such that the sequence \(\left(u(h_{\alpha}(x_n))\right)\) has a limit \(b \in \mathbf{B}\) . Since \(h_{\alpha}\) is continuous, \(h_{\alpha}(a)\) is a limit of the sequence \((h_{\alpha}(x_n))\) in \(\mathbf{E}\) ; therefore by hypothesis \(b = u(h_{\alpha}(a))\) and, since \(\mathrm{F}_{\alpha}\) has the property (GDF), \(u \circ h_{\alpha}\) is continuous.

COROLLARY. --- Every quotient space of a locally convex space having the property (GDF) has the property (GDF).

PROPOSITION 3. --- The strong dual of a reflexive Fréchet space has the property (GDF).

This is a consequence of Prop. 2 and the following lemma (or TVS, IV, §3, No.~4, Prop. 4):

Lemma. --- Let F be a Fréchet space, \(F'\) its strong dual, \(F''\) its bidual. If every subset of \(F''\) , bounded for \(\sigma(F'', F')\) , is contained in the closure (for \(\sigma(F'', F')\) ) of a bounded subset of F, then \(F'\) is the direct limit of a sequence of Banach spaces.

For, let \((\mathrm{V}_{n})\) be a decreasing fundamental sequence of convex, balanced and closed neighborhoods of 0 in F. For every integer n, let \(G_{n}\) be the linear subspace of \(F'\) generated by the polar \(V_{n}^{\circ}\) of \(V_{n}\) . In \(G_{n}\) , \(V_{n}^{\circ}\) is an absorbent convex set, therefore its gauge \(p_{n}\) is a norm on \(G_{n}\) ; moreover, \(V_{n}^{\circ}\) is a complete subset of the strong dual \(F'\) (TVS, III, §3, No.~8, Prop. 11); thus \(G_{n}\) , equipped with the norm \(p_{n}\) , is a Banach space (GT, III, §3, No.~5, Cor. 2 of Prop. 9). We are going to show that the strong topology on \(F'\) is the direct limit of these Banach space topologies on the \(G_{n}\) , or again that, for a strongly closed, balanced, convex subset U of \(F'\) to be a strong neighborhood of 0, it is necessary and sufficient that it absorb each of the \(V_{n}^{\circ}\) . It is obvious that this condition is necessary; to see that it is sufficient, it will suffice to prove that U contains a barrel of \(F'\) . Indeed, its polar \(U^{\circ}\) in \(F''\) will then be bounded for \(\sigma(\mathrm{F}''\mathrm{,}\mathrm{F}')\) , hence, by hypothesis, will be contained in the closure (for \(\sigma(\mathrm{F}''\mathrm{,}\mathrm{F}')\) ) of a bounded subset B of F, from which one can conclude that U (which is closed for \(\sigma(\mathrm{F}^{\prime},\mathrm{F}''\mathrm{)})\) ) contains the strong neighborhood \(B^{\circ}\) of 0 (TVS, II, §6, No.~3, Th. 1 and §8, No.~4).

By hypothesis, for every integer n there exists a number \(\lambda_{n}>0\) such that \(\lambda_{n}V_{n}^{\circ}\subset\frac{1}{2}U\) ; let \(A_{n}\) be the convex envelope of the union of the \(\lambda_{i}V_{i}^{\circ}\) for \(i\leqslant n\) . Then \(A_{n}\subset\frac{1}{2}U\) for every n; let W be the union of the \(A_{n}\) ; W is a convex, balanced, absorbent set contained in \(\frac{1}{2}U\) , and it will suffice to show that its strong closure (which is a barrel) is contained in U.

Thus, let \(x'\) be a point of \(F'\) not belonging to U. Since each of the \(V_n^\circ\) is compact for \(\sigma(F', F)\) , the same is true of the \(A_n\) (TVS, II, §2, No.~6, Prop. 15), and since \(x' \notin 2A_n\) , there exists an element \(x_n\) belonging to the polar of \(A_n\) in F such that \(\langle x', x_n \rangle = 2\) (TVS, II, §5, No.~3, Prop. 4). The sequence \((x_n)\) is bounded in F: for, every \(y' \in F'\) belongs to some \(V_k^\circ\) , consequently \(|\langle y', x_n \rangle| \leqslant \lambda_k^{-1}\) for \(n \geqslant k\) , whence our assertion (TVS, IV, §1, No.~1, Prop. 1). Let C be a bounded, balanced convex set in F containing all the \(x_n\) ; \(C^\circ\) is then a neighborhood of 0 in \(F'\) , and the polar \(C^{\circ\circ}\) of \(C^\circ\) in \(F''\) is compact for \(\sigma(F'', F')\) (TVS, III, §3, No.~4, Cor. 2 of Prop. 4 and No.~5, Prop. 7). Thus one sees that the sequence \((x_n)\) has a cluster point \(x''\) in \(F''\) for \(\sigma(F'', F')\) ; obviously \(\langle x', x'' \rangle = 2\) and, on the other hand, \(x''\) belongs to the polar of \(A_n\) in \(F''\) for every \(n\) , hence to the polar \(W^\circ\) of W in \(F''\) . From this, one concludes that \(x' \notin W^{\circ\circ}\) , hence is not in the closure of W for \(\sigma(F', F'')\) (TVS, II, §6, No.~3, Th. 1 and §8, No.~4), therefore a fortiori for the strong topology, which completes the proof.
% semantic-fragments: legacy-input-seam
\relax{}
